\subsection{Transport in the original coefficient fibre}

We transfer the outer orbit's inequality to the splitting, gap and
saddle-ratio coordinates. The comparisons below may contain $|\nu|$;
the exponential factor absorbs it on the fixed coefficient fibre.

\begin{proposition}[Uniform coefficient and domain transport]
\label{rw:coefficient-transport}
On the ordered cover of Proposition~\ref{rw:ordered-cover},
the fixed original inequality
\eqref{rw:fibre:original-forcing} supplies the relative coefficient
bounds used in the local critical-function estimates, uniformly at
all physical inputs, including
nonclosing inputs. More precisely, for each local center-generator
size $I_{\rm loc}$ used in Sections~7--10 there is a modulus
$\eta_{\rm loc}(\epsilon)\to0$ such that
\[
 |\nu|\le\eta_{\rm loc}(\epsilon)I_{\rm loc}.
\]
The coefficient constants and moduli are common to the finite cover.
The coefficient divisions take place on the larger comparison
domains of Section~5 and the corresponding secondary domains.
The counted inputs remain on the positive-central physical interval.
If a local generator size vanishes, absorption first gives
$\nu=0$ and then the relevant exact center identity.
\end{proposition}
\begin{proof}
Set $\delta_\epsilon=e^{-p_\sharp^2/(20\epsilon)}$.
The original inequality is
\begin{equation}\label{rw:fibre:transport-start}
 |\nu|\le\delta_\epsilon
                 (|\tau_1|+\epsilon|\sigma_1|)
          \le\delta_\epsilon
                 (|\tau_1|+|\sigma_1|).
\end{equation}
We compare the original generators with the local transverse
coefficients in each type of chart.

On regular and strict-focus charts take
$I_{\rm loc}=|\tau_1|+|\sigma_1|$. At a two-lower-root match,
the strict splitting derivative of
Proposition~\ref{rw:finite:two-root-count} makes
$(\bar a,\delta,\nu)$ a bounded inverse chart on a larger box.
Here $\delta$ is the actual strong-intersection splitting.
Single division from its exact reversible axis gives
\[
 |\tau_1|+|\sigma_1|
                 \le C(|\bar a|+|\delta|+|\nu|).
\]
The finite upper saddle uses $(\bar a,\kappa,\nu)$ with the
actual $\kappa=r-1=-\operatorname{tr}/\lambda_+$; the derivative
$\partial_{\bar c}\kappa=-1/\sqrt{\bar m}+O(\epsilon)$ has a
nonzero margin. The same division, already described after
Proposition~\ref{rw:generator-division}, gives
$|\tau_1|+|\sigma_1|\le C(|\bar a|+|\kappa|+|\nu|)$.
Off these center axes the transverse size has a positive lower
bound on each smaller closed chart.

At the finite three-saddle intersection,
\eqref{rw:centers:ideal-comparison} gives
\[
 |\tau_1|+|\sigma_1|+|\nu|
            \asymp |\xi_2|+|\kappa\xi_1|+|\nu|.
\]
Thus take $I_{\rm loc}=|\xi_2|+|\kappa\xi_1|$; the product
retains both intersecting center branches. The double-division
segments and derivative budget are those of Section~5.

At a strict one-infinite simple-root coincidence use the exact
splitting $\delta$ and the actual scaled exponent difference
$g$. On its strict compact row the determinant of this coordinate
change tends to
\[
 \frac{p_c^4}{\bar B(1-\bar Bp_c^2)^3}>0,
\]
up to the positive unit from the regular power coordinate.
The calculation and the equality of the actual leading powers
on the second center are in
Proposition~\ref{rw:infinity:strict-lower-loss}.
Their bounded inverse and single division give
$|\tau_1|+|\sigma_1|\le C(|\delta|+|g|+|\nu|)$.
At the repeated and weak faces use instead the two comparisons
\eqref{rw:centers:single-ideal-comparison}, with local sizes
$|\tau|+|\zeta|$ and $|d_s|+|\varsigma|$, respectively.
Their inverse charts and exact center axes were proved in
Proposition~\ref{rw:center-coordinates}; the repeated comparison
uses \eqref{rw:coordinates:zeta-identity}, rather than the strict-row
determinant. The fixed-mean and lifted-contact comparisons, and the
single weak offset $\bar B+K_*A$ of
Proposition~\ref{rw:domains:weak}, are defined on larger domains
containing every original input and the whole coefficient segments.
At a cubic two-infinite end $\bar B=0,\bar m<0$,
$|\bar B+\bar m|$ has a positive margin. Hence
$|\tau_1|+|\sigma_1|\asymp|\bar a|+|\bar c|$,
the cubic generator size. On other constant-sign pieces
either the same bounded comparison applies or the
transverse size is bounded below; no division at a vanishing
gap is needed.

At the parabolic corner we must retain the scale weights before
comparing generator sizes. Since
$\bar B+\bar m>1/8$, bounded normalized coefficients give
\[
 |\tau_1|+|\sigma_1|
       \le C(|\bar a|+|\bar c|)
       \le C(|\bar a|+|\bar d|+|\nu|).
\]
Inserting this in \eqref{rw:fibre:transport-start} and
absorbing gives
\begin{equation}\label{rw:fibre:weighted-transport}
 |\nu|\le2C\delta_\epsilon(|\bar a|+|\bar d|).
\end{equation}
On the wedge, $\bar d=H\gamma$ with $H\le H_0<1$;
on the secondary charts, $\bar d=\lambda C$ with
$\lambda\le\lambda_*<1$. It follows, uniformly at every
actual input, that the forcing is small relative to the local
sizes $|\bar a|+|\gamma|$ and $|\bar a|+|C|$.
The stronger weighted inequality
\eqref{rw:fibre:weighted-transport} is retained whenever
those local expressions are returned to original coefficients.
The wedge bound in \eqref{rw:secondary:critical-forcing} retains
$|\bar a|+|\bar d|$; the secondary bound uses $|\bar a|+|C|$.
The additional powers in \eqref{rw:secondary:local-forcing} arise
when an actual secondary closing witness supplies its absolute
budget; they are stronger than the bounds used here.

Near a finite secondary upper saddle the exact ratio calculation
\eqref{rw:secondary:saddle-ratio} has a nonzero derivative
in $C$ and yields
\[
 |C|\le C_0\bigl(|\kappa|+
                  \epsilon\lambda|\bar a|
                  +\epsilon\lambda|\nu|\bigr).
\]
Thus the preceding bound implies
$|\nu|\le\widetilde\eta(\epsilon)
 (|\bar a|+|\kappa|)$ after one further absorption.
The estimate therefore uses the actual ratio of the secondary field,
including its coefficient error relative to the limiting linear field.
On the weak and real repeated secondary rows $|C|$ has a
positive margin on the normalized face, so their transverse
size is bounded below. They need no inverse coordinate
division at the vanishing weak eigenvalue. A strict focus
center has its regular division in $(\bar a,C)$, and the
finite saddle center was treated above. These neighborhoods
cover the normalized face with the adjacent strict rows,
so the constants are uniform on their common boundary faces.

In every unweighted comparison we have obtained
$I_{\rm global}\le C(I_{\rm loc}+|\nu|)$.
Choose the common scale with $C\delta_\epsilon<1/2$;
then \eqref{rw:fibre:transport-start} gives
$|\nu|\le2C\delta_\epsilon I_{\rm loc}$.
The weighted comparisons have just given the same conclusion
with their explicitly retained factors. If $I_{\rm loc}=0$,
this inequality gives $\nu=0$ before any division by that
size. The inverse coordinate identity then places the field
on its real center axis, so its actual compared maps agree.
This deals with the zero-size case as well.

These are coefficient inequalities, so they hold also at nonclosing
inputs. All divisions use the larger domains of Section~5, while
saturation supplies positive central heights for the original field.
Taking the largest of the finite coefficient bounds makes the
absorptions and moduli uniform on the whole interval.
\end{proof}
